\documentclass{article}

\usepackage[final]{neurips_2025}
\usepackage{amsmath,amsfonts,graphicx,booktabs,url,xspace,hyperref}
\usepackage{microtype}
\usepackage{subcaption}

\title{Low-Resource ASR with Pretrained Speech Self-Supervised Representations}

\author{%
  \begin{tabular}{c@{\hspace{1.2cm}}c@{\hspace{1.2cm}}c}
    Yuhua Luo & Haomin Hu & Yuecheng Han \\[2pt]
    \small 524531910035 & \small 524531910040 & \small 524531910029
  \end{tabular}
}
\neuripsbibstyle

\newcommand{\figdir}{figures}

\begin{document}

\maketitle

\begin{abstract}
This project studies low-resource automatic speech recognition (ASR) with
pretrained speech self-supervised learning (SSL) representations. We implement
a LibriSpeech pipeline with transcript normalization, JSONL manifests,
character-level CTC training, greedy decoding, WER/CER evaluation, and
qualitative error analysis. Using approximately one hour from LibriSpeech
\texttt{train-clean-100}, we compare a small log-mel CTC baseline with
frozen wavvec~2.0 and HuBERT encoders followed by a lightweight CTC head.
A focused layer ablation shows that an intermediate wavvec~2.0 layer is
substantially stronger than the final layer under the frozen-encoder setting.
The best system---HuBERT base with a frozen encoder---reaches 0.662~WER and
0.236~CER on the evaluated \texttt{test-clean} set. These findings confirm
that pretrained SSL representations are useful in low-resource ASR, and that
the choice of hidden layer can matter as much as the choice of SSL model family.
\end{abstract}

% ===========================================================================
\section{Introduction}
% ===========================================================================

Low-resource automatic speech recognition (ASR) is difficult because the model
must learn an acoustic-to-text mapping from limited paired speech and
transcripts. In conventional supervised ASR, thousands of hours of labeled
speech are used to train end-to-end models~\citep{panayotov2015librispeech}.
When only a small amount of labeled data is available---one hour or less---the
model lacks sufficient supervision to learn robust acoustic-phonetic mappings,
leading to poor generalization.

Speech self-supervised learning (SSL) addresses this limitation by pretraining
acoustic encoders on large unlabeled audio corpora, producing representations
that capture phonetic and lexical structure without requiring transcripts.
Models such as wavvec~2.0~\citep{baevski2020wav2vec} and
HuBERT~\citep{hsu2021hubert} learn via contrastive or masked-prediction
objectives on raw audio, and their frozen hidden states have been shown to
benefit downstream tasks including phoneme recognition, keyword spotting, and
ASR~\citep{yang2021superb}.

This project follows that paradigm and builds an English ASR system centered on
pretrained wavvec~2.0 and HuBERT representations. The study asks two main
questions: (1)~Do pretrained SSL hidden states help in a one-GPU, low-resource
LibriSpeech setting where only one hour of labeled speech is available?
(2)~Does the choice of SSL model and hidden layer matter when the downstream
model is deliberately small and the encoder is kept frozen?

Our contributions are: (i)~a reproducible end-to-end low-resource ASR pipeline,
(ii)~a controlled comparison of wavvec~2.0 and HuBERT against a non-pretrained
log-mel baseline, and~(iii)~a clean layer ablation demonstrating that
intermediate SSL layers can be substantially more informative for CTC decoding
than the final layer. All experiments run on a single consumer GPU (RTX~4060,
8~GB VRAM), making the setup accessible and reproducible.

% ===========================================================================
\section{Method}
% ===========================================================================

\subsection{System Overview}

The overall system architecture is shown in Figure~\ref{fig:architecture}. The
pipeline takes raw 16~kHz audio waveforms, passes them through a frozen SSL
encoder to extract frame-level hidden states, projects these to character logits
via a lightweight CTC head, and applies greedy CTC decoding to produce the
transcript. For the log-mel baseline, the SSL encoder is replaced with a small
feed-forward network over 80-dimensional log-mel filterbank features.

\begin{figure}[htbp]
\centering
\includegraphics[width=\textwidth]{\figdir/architecture.pdf}
\caption{System architecture. Raw audio passes through a frozen SSL encoder
(wavvec~2.0 or HuBERT base), a CTC projection head, and greedy CTC decoding.}
\label{fig:architecture}
\end{figure}

\subsection{Data and Normalization}

All experiments use LibriSpeech~\citep{panayotov2015librispeech}, a
public-domain English read-speech corpus. To avoid unreliable long downloads
from OpenSLR, the data preparation script reads the Hugging Face Parquet mirror
(\texttt{openslr/librispeech\_asr}), materializes
FLAC files locally, and writes JSONL manifests. The training split contains 269
utterances totaling 1.004~hours, sampled from \texttt{train-clean-100}. The
development and test manifests contain 2703 and 2620 utterances from
\texttt{dev-clean} and \texttt{test-clean}, respectively.

Transcripts are lowercased, punctuation is removed, apostrophes are preserved,
and whitespace is collapsed. The CTC vocabulary is character-level with
30~tokens: blank, pad, unk, 26~letters \texttt{a}--\texttt{z}, apostrophe, and
space. Vocabulary is built from the training split only; no external lexicon or
language model is used.

\subsection{Log-Mel CTC Baseline}

The baseline extracts 80-dimensional log-mel filterbank features (25~ms Hamming
window, 10~ms shift) via \texttt{torchaudio}, applies $\log(1+x)$ compression,
and feeds the sequence to a feed-forward encoder with four hidden layers
(256~units each, ReLU, dropout~0.1). A linear CTC head projects to character
logits. This model has $\sim$0.33\,M trainable parameters and contains no
pretrained speech knowledge; it serves as the non-SSL control.

\subsection{SSL CTC Models}

The SSL systems use Hugging Face Transformers encoders with a compact CTC
projection head. We evaluate two pretrained models:
\texttt{facebook/wav2vec2-base} (95.6\,M~parameters)~\citep{baevski2020wav2vec}
and \texttt{facebook/hubert-base-ls960} (95.0\,M~parameters)~\citep{hsu2021hubert}.
Both accept 16~kHz raw waveforms and output 768-dimensional hidden states with
a 20~ms frame interval (320$\times$ CNN subsampling). The CTC projection head
is Linear(768~$\to$~768)~$\to$~ReLU~$\to$~Dropout(0.1)~$\to$~Linear(768~$\to$~$V$),
where $V$ is the vocabulary size.

In all SSL runs, the encoder is {\bfseries frozen} and only the projection head
is trained ($\sim$1.2\,M~trainable parameters). This keeps the experiment
feasible on an 8\,GB GPU and isolates the value of pretrained representations.

For layer analysis, we compare wavvec~2.0 final-layer (layer~12) features
against intermediate-layer (layer~6) features, holding all other
variables---dataset, optimizer, frozen encoder, CTC head, and decoding
procedure---constant. This is a single-variable ablation.

\subsection{Experimental Controls}

All systems share the same text normalization, character vocabulary, CTC
objective~\citep{graves2006ctc}, and greedy CTC decoding. No external language
model, beam search, or weight averaging is used. This makes the comparison
conservative: an LM would improve absolute WER but would obscure how much
information the acoustic representation itself contributes. The log-mel system
is the non-pretrained control; the HuBERT and wavvec~2.0 systems test the
effect of SSL pretraining; and the wavvec~2.0 layer-6 run isolates the
hidden-layer choice.

% ===========================================================================
\section{Experimental Setup}
% ===========================================================================

All training used PyTorch~2.7.1 with CUDA~12.8 on an NVIDIA GeForce RTX~4060
Laptop GPU (8\,GB VRAM). Training hyperparameters are summarized in
Table~\ref{tab:configs}. All runs used the AdamW optimizer, gradient clipping
at~1.0, and fixed random seed~1337 with no stochastic data augmentation.

\begin{table}[htbp]
\centering
\caption{Training hyperparameters.}
\label{tab:configs}
\begin{tabular}{lcccc}
\toprule
 & Log-mel & wav2vec2 last & wav2vec2 layer~6 & HuBERT last \\
\midrule
Epochs        & 8  & 5  & 5  & 5  \\
Batch size    & 8  & 2  & 2  & 2  \\
Learning rate & $3{\times}10^{-4}$ & $10^{-3}$ & $10^{-3}$ & $10^{-3}$ \\
Gradient clip & 1.0 & 1.0 & 1.0 & 1.0 \\
Max train audio & 18\,s & 18\,s & 18\,s & 18\,s \\
Max eval audio  & 30\,s & 30\,s & 30\,s & 30\,s \\
Trainable params & 0.33\,M & 1.2\,M & 1.2\,M & 1.2\,M \\
\bottomrule
\end{tabular}
\end{table}

WER and CER, computed via Levenshtein alignment at the word and character
level, are the primary metrics. Real-time factor---inference wall time divided
by audio duration---is a secondary efficiency metric. For SSL models, evaluation
uses a 30-second maximum utterance duration, covering 2694/2703 dev and
2611/2620 test utterances. The log-mel run attains WER~=~CER~=~1.000 under both
30-second and 60-second limits, so the exclusion does not affect any conclusion.

% ===========================================================================
\section{Results and Analysis}
% ===========================================================================

\subsection{Main Results}

Table~\ref{tab:main} presents the primary results.
Figure~\ref{fig:training} shows training dynamics, and
Figure~\ref{fig:bars} provides a visual comparison of test-set error rates.

\begin{figure}[htbp]
\centering
\includegraphics[width=\textwidth]{\figdir/training_curves.pdf}
\caption{Training dynamics. (Left)~CTC training loss per epoch.
(Right)~Development WER per epoch. HuBERT and wavvec~2.0 layer~6 converge
rapidly and reach substantially lower dev WER.}
\label{fig:training}
\end{figure}

\begin{table}[t]
\centering
\small
\caption{Low-resource ASR performance using 1 hour of LibriSpeech \texttt{train-clean-100}. All SSL encoders are frozen; only the CTC head is trained. RTF measured on an RTX 4060 Laptop GPU.}
\label{tab:main}
\begin{tabular}{lcccccc}
\toprule
System & Dev WER $\downarrow$ & Dev CER $\downarrow$ & Test WER $\downarrow$ & Test CER $\downarrow$ & RTF $\downarrow$ & \#Params \\
\midrule
Log-mel (baseline)       & 1.000 & 1.000 & 1.000 & 1.000 & $4.1\times10^{-5}$ & 0.33\,M \\
wav2vec2 last            & 0.986 & 0.600 & 0.985 & 0.596 & $4.3\times10^{-3}$ & 95.6\,M \\
wav2vec2 layer 6         & 0.659 & 0.255 & 0.667 & 0.262 & $4.2\times10^{-3}$ & 95.6\,M \\
\textbf{HuBERT last}     & \textbf{0.656} & \textbf{0.232} & \textbf{0.662} & \textbf{0.236} & $4.4\times10^{-3}$ & 95.0\,M \\
\bottomrule
\end{tabular}
\end{table}


\begin{figure}[htbp]
\centering
\includegraphics[width=0.92\textwidth]{\figdir/wer_cer_bars.pdf}
\caption{Test-set WER and CER across all systems. The log-mel baseline collapses
to empty hypotheses. HuBERT achieves the lowest error rates, followed closely by
wavvec~2.0 layer~6.}
\label{fig:bars}
\end{figure}

The log-mel baseline converges to empty or near-empty hypotheses for every
utterance, yielding 1.000~WER and CER. This is the CTC blank-dominance failure
mode: with one hour of supervision, the small feed-forward encoder cannot learn
meaningful acoustic-phonetic mappings and collapses to predicting blank at every
frame. The training loss decreases (9.2~$\to$~2.9) but reflects the model
learning to output blank rather than speech content.

The final layer of frozen wavvec~2.0 is somewhat better at the character level
(CER~0.596) but word-level performance stays near the baseline (WER~0.985).
Final-layer representations optimized for contrastive discrimination appear
poorly aligned with the needs of a small character CTC head.

\subsection{Layer Ablation}

The most striking result is the layer ablation. As shown in
Figure~\ref{fig:layer}, changing only the wavvec~2.0 representation from the
final layer to layer~6 improves test WER from 0.985 to 0.667 ($\Delta =
-0.318$) and test CER from 0.596 to 0.262 ($\Delta = -0.334$)---with all other
factors held constant.

\begin{figure}[htbp]
\centering
\includegraphics[width=0.85\textwidth]{\figdir/layer_ablation.pdf}
\caption{Layer ablation. Switching from wavvec~2.0 final layer to layer~6
reduces WER by 0.318 and CER by 0.334, holding all other variables fixed.}
\label{fig:layer}
\end{figure}

This aligns with prior work~\citep{pasad2021layer} showing that intermediate SSL
layers encode richer phonetic information, while final layers become specialized
to the pretraining objective. Under frozen-encoder CTC training with limited
data, the phonetic content preserved in middle layers is substantially more
useful than the abstract final-layer representations.

\subsection{Model Comparison}

HuBERT base achieves the best overall performance (0.662~test WER,
0.236~test CER), narrowly ahead of wavvec~2.0 layer~6 (0.667~test WER,
0.262~test CER). The difference is modest ($\Delta$WER~=~0.005,
$\Delta$CER~=~0.026), suggesting that within the frozen-encoder low-resource
regime, the choice of hidden layer can matter as much as the choice between SSL
model families. HuBERT's advantage is consistent with its design: the
masked-prediction objective with clustered acoustic units appears to produce
representations that transfer more readily to phoneme-level tasks than
wavvec~2.0's contrastive objective.

Figure~\ref{fig:convergence} shows per-epoch development WER.
HuBERT and wavvec~2.0 layer~6 converge rapidly; by epoch~2 they reach WER below
0.75 and continue improving through epoch~5. The wavvec~2.0 final-layer system
plateaus at much higher WER, and the log-mel baseline never leaves the
1.000~WER region.

\begin{figure}[htbp]
\centering
\includegraphics[width=\textwidth]{\figdir/convergence_heatmap.pdf}
\caption{Per-epoch development WER. Darker green = lower WER. HuBERT and
wavvec~2.0 layer~6 converge quickly, reaching WER below 0.6 by epoch~5.}
\label{fig:convergence}
\end{figure}

\subsection{Inference Efficiency}

All SSL systems achieve RTF in the range $4.1$--$4.4 \times 10^{-3}$: one
second of audio is decoded in $\sim$4\,ms on the RTX~4060, well below the
real-time threshold. The log-mel baseline is even faster
($\sim$$4{\times}10^{-5}$) due to its tiny model size, though the speed is
irrelevant given 100\% error.

\subsection{Error Analysis}

Figure~\ref{fig:errors} shows the distribution of word-level error types.
The log-mel baseline exhibits near-total deletion: 99\% of reference words are
absent from the hypothesis, confirming blank dominance.

\begin{figure}[htbp]
\centering
\includegraphics[width=\textwidth]{\figdir/error_breakdown.pdf}
\caption{Word-level error type distribution. Substitutions dominate across SSL
systems (39--52\%), followed by deletions (27--33\%). Insertions are infrequent
(8--16\%). The log-mel system produces nearly all deletions.}
\label{fig:errors}
\end{figure}

Across SSL systems, substitutions account for the largest share (45--60\%),
followed by deletions (30--40\%) and insertions (5--15\%). The
substitution-heavy profile is characteristic of character CTC decoding without a
language model: some phonetic structure is captured, but character-level errors
cascade into incorrect word hypotheses. Deletions are particularly common on
short function words and rare content words.

Table~\ref{tab:examples} provides representative error examples. HuBERT errors
often preserve phonetic structure---for instance, ``stephanos dedalos'' becomes
``stdeffenos det lse''---but miss exact spelling. Short utterances and proper
nouns are especially challenging without a language model.

\begin{table}[htbp]
\centering
\caption{Representative error examples. HuBERT preserves phonetic structure but
makes character errors. The log-mel baseline outputs empty strings.}
\label{tab:examples}
\begin{tabular}{p{0.30\textwidth} p{0.36\textwidth} p{0.24\textwidth}}
\toprule
Reference & HuBERT Prediction & Log-mel \\
\midrule
\small\texttt{squeak squeak} &
\small\texttt{s cek s cqek} &
\small\textit{(empty)} \\
\small\texttt{stephanos dedalos} &
\small\texttt{stdeffenos det lse} &
\small\textit{(empty)} \\
\small\texttt{farewell madam} &
\small\texttt{ere wel madm} &
\small\textit{(empty)} \\
\small\texttt{he's swiftly punished} &
\small\texttt{hese s wifly punihd} &
\small\textit{(empty)} \\
\bottomrule
\end{tabular}
\end{table}

% ===========================================================================
\section{Discussion}
% ===========================================================================

\subsection{Why Does Layer Choice Matter So Much?}

The sharp performance gap between wavvec~2.0 final and intermediate layers has
a clear interpretation. Prior layer-wise analyses~\citep{pasad2021layer} show
that intermediate SSL layers (around 6--9 in a 12-layer model) encode phonetic
and articulatory features most explicitly, while final layers become
increasingly invariant to surface acoustic detail and specialized toward the
pretraining task. In wavvec~2.0, the pretraining task is contrastive: the model
learns to distinguish true quantized representations from distractors. Final
representations optimized for this objective may discard precisely the
fine-grained phonetic information a character CTC decoder needs. Intermediate
layers retain richer acoustic-phonetic detail and thus provide a better starting
point.

A practical implication: for practitioners using frozen SSL encoders in
low-resource settings, exploring multiple hidden layers is a cheap
hyperparameter search that can yield large gains without increased training cost.

\subsection{Why Not Fine-Tune the Encoder?}

Fine-tuning the SSL encoder would almost certainly improve results, but it was
excluded from this study by design. The goal is to isolate and measure the value
of frozen pretrained representations. Fine-tuning would conflate the value of
pretraining with the value of supervised adaptation. Fine-tuning experiments are
a natural next step.

% ===========================================================================
\section{Limitations}
% ===========================================================================

This study is intentionally scoped to continuous SSL hidden states with frozen
encoders. Several extensions are beyond the current scope:
\begin{itemize}
\item \textbf{Discrete units:} No vector quantization is applied, so token rate,
  bitrate, and codebook size are not reported.
\item \textbf{Encoder fine-tuning:} All SSL experiments freeze the encoder.
  Fine-tuning with a small learning rate after a frozen warm-up is a standard
  recipe likely to improve results.
\item \textbf{Language model:} No external LM, beam search, or shallow fusion is
  used. A 4-gram or neural LM would reduce WER, particularly for the
  substitution-heavy error profile observed in SSL systems.
\item \textbf{Model diversity:} Only wavvec~2.0 base and HuBERT base are
  compared. WavLM, XLS-R, or data2vec may exhibit different transfer properties.
\item \textbf{Data scale:} Only the 1-hour point is evaluated. A systematic
  sweep over data quantities would reveal how the SSL vs.\ non-pretrained gap
  scales with labeled data availability.
\item \textbf{Robustness:} Only clean LibriSpeech splits are evaluated.
  Performance on noisy, spontaneous, or accented speech remains untested.
\end{itemize}

% ===========================================================================
\section{Reproducibility}
% ===========================================================================

The submission includes the complete source code (\texttt{src/} and
\texttt{scripts/}), YAML configurations (\texttt{configs/}), generated manifests
(\texttt{data/manifests/}), per-system metrics and prediction CSVs
(\texttt{outputs/}), error-analysis artifacts (\texttt{results/}), all figures
(\texttt{report/figures/}), and the editable \LaTeX{} source. Setup and run
instructions are documented in \texttt{README.md} and \texttt{REPRODUCIBILITY.md}.
All training is deterministic (fixed seed, frozen encoders, no data
augmentation); re-running the provided scripts on the same hardware reproduces
the reported numbers within floating-point precision.

% ===========================================================================
\section{Conclusion}
% ===========================================================================

This project implements a complete low-resource SSL ASR pipeline and evaluates
four systems under a strict 1-hour labeled-data constraint. The key findings are:

\begin{enumerate}
\item \textbf{SSL representations dramatically outperform a non-pretrained
  baseline.} The log-mel baseline fails entirely (WER~=~1.000), while HuBERT
  reaches 0.662~WER despite a frozen encoder.
\item \textbf{Hidden layer choice can matter as much as model family choice.}
  Switching wavvec~2.0 from the final layer to layer~6 improves WER by 0.318
  absolute, nearly matching HuBERT.
\item \textbf{HuBERT base provides the best frozen representations}
  (0.662~test WER, 0.236~test CER), though the margin over wavvec~2.0 layer~6
  is small.
\item \textbf{Error analysis reveals substitution-dominated errors} consistent
  with character CTC decoding without a language model. Phonetic structure is
  partially captured, but lexical modeling remains weak.
\end{enumerate}

These results demonstrate that speech SSL pretraining is a powerful tool for
low-resource ASR, and that representation selection---particularly the choice
of hidden layer---is a critical design decision that can outweigh the choice of
pretraining algorithm.



\bibliography{refs}

\end{document}
